\newpage

\subsection{Artificial Intelligence Application}

The emergence of Artificial Intelligence (AI) has enabled more efficient workflows by reducing the need for human involvement in tedious and repetitive tasks. At the same time, there has been a growing reliance on AI technologies. The author acknowledges that AI tools were utilized throughout this project under their full control and responsibility.

The primary use of AI was as a rapid research aid during the implementation phase. When new technical challenges arose, AI tools were employed to assist in identifying potential solutions and diagnosing issues. This facilitated faster bug resolution by guiding the author toward the root causes of problems.

A second key application of AI was in the proofreading and refinement of the thesis documentation. Given the extensive amount of content required within a limited timeframe, AI tools were used to detect grammatical errors and improve the overall clarity and coherence of the text.

These represent the main areas in which AI was applied. However, its use was not strictly limited to them. AI also supported the author in learning and understanding unfamiliar concepts, always complemented by manual verification through reliable sources. 

Additionally, AI-based coding assistants, such as GitHub Copilot, were used during development to automate the generation of repetitive code structures, thereby improving productivity. Despite this assistance, all AI-generated code was critically evaluated and validated by the author to ensure correctness and alignment with project requirements.
